\documentclass{article}
\usepackage[a4paper]{geometry}
\usepackage{fancyhdr}
\pagestyle{fancy}
\lhead{Homologien}
\rhead{Februar 2026}
\begin{document}
\section{Homologien}
\emph{Homologien} beschreiben ähnlichkeiten in Strukturen der (oftmals erwachsenen) Lebewesen unterschiedlicher Arten welche genutzt werden können um Verwandschaften zu finden und zu begründen. Sie entstehen durch sog. \emph{divergente} Entwicklung; hier wird aus einem gemeinsamen Vorfahren mit dieser ähnlichen Struktur nun unterschiedliche Arten.
 
\subsection{Homologiekriterien} 
Es kann mit der Lage, der spezifischen Qualität und der Stetigkeit argumentiert werden.
\begin{description}
 \item[Lage] Strukturen sind wahrscheinlich homolog, wenn sie in Körper relativ zueinander Positioniert sind.
 \item[Spezifische Qualität] Strukturen sind wahrscheinlich homolog, wenn sie in viele unterschiedliche einzelne Merkmale aufgeteilt werden können, welche sich gegenseitig ähnlich sind.
 \item[Stetigkeit] Auch wenn die ersten beiden Kriterien nicht erfüllt wurden können Strukturen homolog sein wenn sie sich evolutiv verbinden lassen; wenn eine Übergangsform vom einen Strukturaufbau zum anderen gefunden wird.
\end{description} 
 
\subsection{Analogien}
Analogien scheinen Homologien zu sein, beschreiben aber das Gegenteil. Sie sind die Ähnlichkeiten und Gemeinsamkeiten die durch \emph{Konvergenz} entstanden sind.
 
Hier gibt es nicht einen Vorfahren welcher die bestimmte Merkmal entwickelt hatte sondern unterschiedliche Vorfahren, welche unabhängig voneinander ein ähnliches Merkmal entwickelt hatten, weil auf beide Arten ein ähnlicher Selektionsdruck herrschte.  
\end{document}